% SEGMENT OLP-0607-S01
% Part: methods
% Chapter: proofs
% Section: example-1

\documentclass[../../../include/open-logic-section]{subfiles}

\begin{document}

\olfileid{mth}{prf}{ex1}

\olsection{ஓர் எடுத்துக்காட்டு}

முதல் எடுத்துக்காட்டு கணங்களின் ஒன்றிப்பும் வெட்டும் பற்றிய
பின்வரும் எளிய உண்மை. வரையறைகளை விரித்துரைத்தல்,
இணைப்புக் கூற்றுகளையும் அனைத்தையும் பற்றிய கூற்றுகளையும்
நிறுவுதல், நிகழ்வுகளைப் பிரித்து நிறுவுதல் ஆகியவற்றை
இது விளக்கும்.

\begin{prop}
எந்தக் கணங்கள் $A$, $B$, $C$ க்கும்,
$A \cup (B \cap C) = (A \cup B) \cap (A \cup C)$.
\end{prop}

இதை நிறுவுவோம்!{}

% SEGMENT OLP-0607-S02
\begin{proof}
எந்தக் கணங்கள் $A$, $B$, $C$ க்கும்,
$A \cup (B \cap C) = (A \cup B) \cap (A \cup C)$
என்று காட்ட வேண்டும்.
\begin{quote}
முதலில் முன்மொழிவிலுள்ள ``$=$'' குறியின் வரையறையை
விரிக்கிறோம். இரு கணங்கள் ஒன்றே என்று நிறுவுவது,
அவற்றுக்கு ஒரே !!{element}s உள்ளன என்று காட்டுவதே.
அதாவது $A \cup (B \cap C)$ இன் எல்லா
!!{element}s உம் $(A \cup B) \cap (A \cup C)$
இன் !!{element}s ஆகவும், மறுதிசையிலும் இருக்க வேண்டும்.
``மறுதிசையிலும்'' என்பதால் $(A \cup B) \cap (A \cup C)$
இன் ஒவ்வோர் !!{element} உம் $A \cup (B \cap C)$
இன் !!a{element} ஆக வேண்டும். வரையறையை விரிக்கும்போது
ஓர் இணைப்புக் கூற்றை நிறுவ வேண்டியிருப்பது தெரிகிறது.
அதைப் பதிவுசெய்வோம்:
\end{quote}
வரையறையின்படி, $A \cup (B \cap C) = (A \cup B) \cap (A \cup C)$
எனில், எனில் மட்டுமே, $A \cup (B \cap C)$ இன் ஒவ்வோர்
!!{element} உம் $(A \cup B) \cap (A \cup C)$ இன்
!!a{element} ஆகவும், $(A \cup B) \cap (A \cup C)$
இன் ஒவ்வோர் !!{element} உம் $A \cup (B \cap C)$
இன் !!a{element} ஆகவும் இருக்கும்.
\begin{quote}
இது இணைப்புக் கூற்று என்பதால், அதன் ஒவ்வொரு கூறையும்
தனித்தனியாக நிறுவ வேண்டும். முதலில் $A \cup (B \cap C)$
இன் ஒவ்வோர் !!{element} உம் $(A \cup B) \cap (A \cup C)$
இன் !!a{element} என்று நிறுவுவோம்.

இது அனைத்தையும் பற்றிய கூற்று. எனவே $A \cup (B \cap C)$
இன் தன்னிச்சையான ஓர் !!{element} ஐ எடுத்து, அது
$(A \cup B) \cap (A \cup C)$ இன் !!a{element}
ஆக வேண்டும் என்று காட்டுவோம். அந்தத் தன்னிச்சையான
!!{element} ஐ $z$ என்று பெயரிடுவோம். நிறுவல்
இப்படித் தொடர்கிறது:
\end{quote}

% SEGMENT OLP-0607-S03
முதலில் $A \cup (B \cap C)$ இன் ஒவ்வோர்
!!{element} உம் $(A \cup B) \cap (A \cup C)$
இன் !!a{element} என்று நிறுவுகிறோம்.
$z \in A \cup (B \cap C)$ என்க. $z \in (A \cup B) \cap (A \cup C)$
என்று காட்ட வேண்டும்.
\begin{quote}
இப்போது $\cup$, $\cap$ ஆகியவற்றின் வரையறைகளை
விரிக்க வேண்டும். எடுத்துக்காட்டாக,
$\cup$ இன் வரையறை:
$A \cup B = \Setabs{z}{z \in A \text{ அல்லது } z \in B}$.
``$A \cup (B \cap C)$'' க்கு இதைப் பயன்படுத்தும்போது,
வரையறையில் ``$B$'' வகித்த இடத்தை ``$B \cap C$''
வகிக்கிறது. ஆகவே
$A \cup (B \cap C) = \Setabs{z}{z \in A \text{ அல்லது } z \in B \cap C}$.
எனவே $z \in A \cup (B \cap C)$ என்ற எடுகோளின்
பொருள் $z \in \Setabs{z}{z \in A \text{ அல்லது } z \in B \cap C}$.
மேலும் $z \in \Setabs{z}{\dots z\dots}$ எனில்,
எனில் மட்டுமே, \dots $z$ \dots; இங்கு அது
$z \in A$ அல்லது $z \in B \cap C$ என்பதாகும்.
\end{quote}
$\cup$ இன் வரையறையின்படி, $z \in A$ அல்லது
$z \in B \cap C$.
\begin{quote}
இது ஒரு பிரிப்புக் கூற்று. எனவே நிகழ்வுகளைப்
பிரித்து நிறுவுவது பயனுள்ளது. இரு நிகழ்வுகளிலும்
நாம் நோக்கும் முடிவு, அதாவது
``$z \in (A \cup B) \cap (A \cup C)$'', கிடைக்கிறது
என்று காட்டுவோம்.
\end{quote}

% SEGMENT OLP-0607-S04
நிகழ்வு 1: $z \in A$ என்க.
\begin{quote}
எடுகோள்களிலிருந்து மட்டும் மேற்கொண்டு பெரிதாகச்
செய்ய முடியாது. எனவே நோக்கிய முடிவைப் பார்ப்போம்.
$z \in (A \cup B) \cap (A \cup C)$ என்று காட்ட வேண்டும்.
$\cap$ இன் வரையறையின்படி, $z \in (A \cup B) \cap (A \cup C)$
என்று காட்ட வேண்டுமானால், அது $(A \cup B)$
இலும் $(A \cup C)$ இலும் இருக்க வேண்டும்.
ஆனால் $z \in A \cup B$ எனில், எனில் மட்டுமே,
$z \in A$ அல்லது $z \in B$; நிகழ்வு~1 இன்
எடுகோளாகவே $z \in A$ உள்ளது. அதே காரணத்தால்,
$B$ க்குப் பதில் $C$ ஐ வைத்தால் $z \in A \cup C$.
இந்தச் சிந்தனை வேண்டிய முடிவிலிருந்து பின்னோக்கிச்
சென்றது; இப்போது நிறுவலுக்குத் தேவையான முன்னோக்கிய
திசையில் எழுதுவோம்.
\end{quote}
$z \in A$ என்பதால், $z \in A$ அல்லது $z \in B$;
ஆகவே $\cup$ இன் வரையறையின்படி $z \in A \cup B$.
அதேபோல் $z \in A \cup C$. எனவே $\cap$ இன்
வரையறையின்படி $z \in (A \cup B) \cap (A \cup C)$.
\begin{quote}
நிகழ்வுகளைப் பிரிக்கும் நிறுவலின் முதல் நிகழ்வு முடிந்தது.
இப்போது $z \in B \cap C$ என்ற இரண்டாவது நிகழ்விலும்
அதே முடிவைப் பெற வேண்டும்.
\end{quote}

% SEGMENT OLP-0607-S05
நிகழ்வு 2: $z \in B \cap C$ என்க.
\begin{quote}
மீண்டும் இரு கணங்களின் வெட்டைக் கையாள்கிறோம்.
$\cap$ இன் வரையறையைப் பயன்படுத்துவோம்:
\end{quote}
$z \in B \cap C$ என்பதால், $\cap$ இன்
வரையறையின்படி $z$ ஆனது $B$ மற்றும் $C$ இரண்டின்
!!a{element} ஆக வேண்டும்.
\begin{quote}
மீண்டும் முடிவைப் பார்ப்போம். $z$ ஆனது
$(A \cup B)$ இலும் $(A \cup C)$ இலும்
இருக்கிறது என்று காட்ட வேண்டும். இம்முறையும்
விடை உடனே கிடைக்கிறது.
\end{quote}
$z \in B$ என்பதால் $z \in (A \cup B)$.
$z \in C$ என்பதால் $z \in (A \cup C)$ உம்.
ஆகவே $z \in (A \cup B) \cap (A \cup C)$.
\begin{quote}
இங்கும் $\cup$, $\cap$ இன் வரையறைகளைப் பயன்படுத்தினோம்.
ஆனால் அவற்றை ஏற்கெனவே நினைவுபடுத்தியுள்ளோம்; $z$
இரு கணங்களில் ஒன்றில் இருந்தால் அவற்றின் ஒன்றிப்பிலும்
இருக்கும் என்று ஏற்கெனவே காட்டினோம். ஆகவே செய்ததை
மீண்டும் இவ்வளவு வெளிப்படையாக விளக்கத் தேவையில்லை.

இரண்டாவது நிகழ்வும் முடிந்தது. எனவே முதல் முடிவை
இப்போது கூறலாம்.
\end{quote}
ஆக $z \in A \cup (B \cap C)$ எனில்,
$z \in (A \cup B) \cap (A \cup C)$.
\begin{quote}
இப்போது மறுதிசையை மட்டும் காட்ட வேண்டும்:
$(A \cup B) \cap (A \cup C)$ இன் ஒவ்வோர்
!!{element} உம் $A \cup (B \cap C)$ இன்
!!a{element} ஆகும். முன்புபோல, முதல் கணத்தின்
தன்னிச்சையான ஓர் !!{element} ஐ எடுத்துக்கொண்டு,
அது இரண்டாவது கணத்திலும் உள்ளது என்று காட்டி
இந்த அனைத்தையும் பற்றிய கூற்றை நிறுவுவோம்.
செய்யவிருப்பதைக் கூறுவோம்.
\end{quote}

% SEGMENT OLP-0607-S06
இப்போது $z \in (A \cup B) \cap (A \cup C)$ என்க.
$z \in A \cup (B \cap C)$ என்று காட்ட வேண்டும்.
\begin{quote}
இப்போது $z \in (A \cup B) \cap (A \cup C)$
என்ற எடுகோளிலிருந்து தொடங்குகிறோம். நிறுவலின் முதல்
பகுதியில் பயன்படுத்திய அதே $z$ ஐ இங்கும் பயன்படுத்துவது
குழப்பமாக இருக்க வேண்டியதில்லை. அப்பகுதியை முடித்தவுடன்
அங்கு எடுத்துக்கொண்ட எல்லா எடுகோள்களும் முடிவுற்றன;
இப்போது $z$ பற்றி புதிய எடுகோள்களை எடுக்கலாம்.
இது குழப்பமாக இருந்தால், தொடரும் பகுதியில் $z$ க்குப்
பதிலாக வேறு மாறியைப் பயன்படுத்துங்கள்.

$\cap$ இன் வரையறையின்படி $z$ ஆனது $A \cup B$
இலும் $A \cup C$ இலும் உள்ளது. $\cup$ இன்
வரையறையையும் விரித்தால், $z \in A$ அல்லது
$z \in B$; மேலும் $z \in A$ அல்லது $z \in C$.
மீண்டும் நிகழ்வுகளைப் பிரிக்க வேண்டியதுபோல் தெரிகிறது;
ஆனால் ``மேலும்'' என்ற இணைப்பு குழப்பம் தரலாம்.
$z$ ஆனது $A$, $B$, $C$ ஆகியவற்றில் ஏதோ ஒன்றில்
இருக்கிறது என்பதே இதன் பொருள் என்று எண்ணலாம்.
அது தவறு. கவனமாக, ஒவ்வொரு பிரிப்புக் கூற்றையும்
தனித்தனியாகப் பார்ப்போம்.
\end{quote}
$\cap$ இன் வரையறையின்படி $z \in A \cup B$
மற்றும் $z \in A \cup C$. $\cup$ இன் வரையறையின்படி
$z \in A$ அல்லது $z \in B$. நிகழ்வுகளைப்
பிரிக்கிறோம்.
\begin{quote}
முதல் பிரிப்புக் கூற்றில் கவனம் செலுத்துவதால்,
$A \cup C$ ஐ விரித்து இரண்டாவது பிரிப்புக் கூற்றை
இன்னும் எழுதவில்லை. இப்போது அது தேவையுமில்லை.
முதல் நிகழ்வு $z \in A$. ஒரு கணத்தின்
!!a{element}, அந்தக் கணத்தையும் வேறு எந்தக்
கணத்தையும் ஒன்றிப்பதால் வரும் கணத்தின்
!!a{element} ஆகும். எனவே நிகழ்வு~1 எளிது:
\end{quote}
நிகழ்வு 1: $z \in A$ என்க. அப்போது
$z \in A \cup (B \cap C)$ என்பது பின்தொடர்கிறது.

% SEGMENT OLP-0607-S07
\begin{quote}
இப்போது இரண்டாவது நிகழ்வு $z \in B$.
இங்கு இரண்டாவது $\cup$ ஐ விரித்து, மீண்டும்
நிகழ்வுகளைப் பிரிப்போம்:
\end{quote}
நிகழ்வு 2: $z \in B$ என்க. $z \in A \cup C$
என்பதால், $z \in A$ அல்லது $z \in C$.
மேலும் நிகழ்வுகளைப் பிரிக்கிறோம்:

நிகழ்வு 2a: $z \in A$. மீண்டும்
$z \in A \cup (B \cap C)$.
\begin{quote}
இது சற்று வியப்பாக இருக்கலாம். இந்தத் துணைநிகழ்வில்
$z \in B$ என்ற எடுகோள் தேவையே இல்லை;
அது பிரச்சினையல்ல.
\end{quote}
நிகழ்வு 2b: $z \in C$. அப்போது $z \in B$
மற்றும் $z \in C$; எனவே $z \in B \cap C$;
இதனால் $z \in A \cup (B \cap C)$.
\begin{quote}
இரு நிலைகளிலும் நிகழ்வுகளைப் பிரித்த நிறுவல்
முடிந்தது; மறுதிசையும் நிறுவப்பட்டது.
\end{quote}
ஆக $z \in (A \cup B) \cap (A \cup C)$ எனில்,
$z \in A \cup (B \cap C)$.
\end{proof}

\end{document}
